% STATUS: COMPLETE DRAFT
% Appendix C : English / Italian glossary and references (only real, verifiable sources)
\bichapter{Glossary and references}{Glossario e bibliografia}\label{app:glossary}

\section*{Glossary \textperiodcentered\ {\color{itgreen}\itshape Glossario}}
{\small
\renewcommand{\arraystretch}{1.12}
\rowcolors{2}{lightblue!45}{white}
\begin{tabularx}{\textwidth}{@{}>{\raggedright\arraybackslash}X >{\raggedright\arraybackslash\itshape\leavevmode\color{itgreen}}X | >{\raggedright\arraybackslash}X >{\raggedright\arraybackslash\itshape\leavevmode\color{itgreen}}X@{}}
\toprule
\rowcolor{white}
\textbf{English} & \textbf{\upshape Italiano} & \textbf{English} & \textbf{\upshape Italiano}\\
\midrule
ray                     & semiretta                  & triangle                  & triangolo\\
vertex                  & vertice                    & equilateral, isosceles, scalene & equilatero, isoscele, scaleno\\
side (of an angle)      & lato                       & acute, right, obtuse triangle & acutangolo, rettangolo, ottusangolo\\
measure of an angle     & ampiezza                   & legs and hypotenuse       & cateti e ipotenusa\\
degree, minute, second  & grado, primo, secondo      & legs and base (isosceles) & lati obliqui e base\\
protractor              & goniometro                 & vertex angle, base angles & angolo al vertice, angoli alla base\\
convex, reflex angle    & angolo convesso, concavo   & altitude, foot            & altezza, piede\\
zero, full angle        & angolo nullo, angolo giro  & orthocentre               & ortocentro\\
straight angle          & angolo piatto              & median, centroid          & mediana, baricentro\\
adjacent angles         & angoli consecutivi, adiacenti & bisector               & bisettrice\\
vertical angles         & angoli opposti al vertice  & midpoint                  & punto medio\\
complementary           & complementari              & exterior angle            & angolo esterno\\
supplementary           & supplementari              & congruent                 & congruente\\
explementary            & esplementari               & rigid motion              & movimento rigido\\
parallel, perpendicular & parallele, perpendicolari  & slide, turn, flip         & traslazione, rotazione, ribaltamento\\
transversal             & trasversale, secante       & corresponding parts       & elementi omologhi\\
corresponding angles    & angoli corrispondenti      & criterion (criteria)      & criterio (criteri)\\
alternate interior      & alterni interni            & given, to prove           & ipotesi, tesi\\
alternate exterior      & alterni esterni            & proof, QED                & dimostrazione, c.v.d.\\
same side interior      & coniugati interni          & common side               & lato in comune\\
pie chart               & areogramma                 & extension of a side       & prolungamento di un lato\\
\bottomrule
\end{tabularx}
}

\section*{References \textperiodcentered\ {\color{itgreen}\itshape Bibliografia}}
\begin{bi}
The history in this booklet is kept to facts that these works document. Where a claim is only a common hypothesis (the link between 360 and the days of the year), the text says so.
\IT
La storia in questo quaderno si limita ai fatti documentati da queste opere. Dove un'affermazione è solo un'ipotesi diffusa (il legame tra 360 e i giorni dell'anno), il testo lo dice.
\end{bi}
{\small
\begin{enumerate}[label={[\arabic*]}, leftmargin=2.2em, itemsep=2pt]
\item Euclid. \emph{The Thirteen Books of Euclid's Elements}. Translated with commentary by T.\,L.~Heath, 2nd ed., 3 vols. Cambridge University Press, 1926; reprinted by Dover, New York, 1956. (Book I: Prop.~4, SAS; Prop.~5, isosceles triangle; Prop.~8, SSS; Prop.~26, ASA; Prop.~29 and 32, parallels and the angle sum.)
\item Proclus. \emph{A Commentary on the First Book of Euclid's Elements}. Translated by G.\,R.~Morrow. Princeton University Press, 1970. (Source of the ``royal road'' anecdote.)
\item D.~Hilbert. \emph{Grundlagen der Geometrie}. Teubner, Leipzig, 1899. (Takes a form of SAS as an axiom of congruence.)
\item O.~Neugebauer. \emph{The Exact Sciences in Antiquity}. 2nd ed., Brown University Press, Providence, 1957; reprinted by Dover, New York, 1969. (Babylonian sexagesimal arithmetic and its passage into Greek astronomy.)
\item E.~Robson. ``Words and pictures: new light on Plimpton 322''. \emph{The American Mathematical Monthly} 109(2), 2002, pp.~105 to 120. (Dating and context of an Old Babylonian base 60 tablet.)
\item G.\,J.~Toomer (translator). \emph{Ptolemy's Almagest}. Duckworth, London, 1984. (Degrees divided into sexagesimal minutes and seconds.)
\end{enumerate}
}
